% arara: pdflatex
% arara: pdflatex
\documentclass[a4paper,11pt]{pressrelease}

\usepackage[british]{babel}
\usepackage{graphicx}

\PRset{head=below,logo=below right,smashlogo,ruled=false}

\PRlogo{\includegraphics[height=2cm]{example-image}}

\PRcompany{Some Company}
\PRdepartment{Some Department}
\PRcontact{Ann Other}
\PRlocation{Some City}
\PRaddress{1 The Street\\The Town\\AB1 2YZ}
\PRphone{01234 56789}
\PRmobile{07123456789}
\PRfax{01234 56788}
\PRurl{http://www.some-company.com/}
\PRemail{ann.other@some-company.com}
\PRhours{9:00--17:30 Mon--Fri}

\PRheadline{Some amazing news}
\PRsubheadline{subheading}

\begin{document}

\begin{pressrelease}

This is an example press release. Keep it short and use the third
person. Avoid the use of exclamation marks and all-caps. Put all the
pertinent details in the first paragraph. Answer who, what, when,
where and why. This sample file uses:
\begin{verbatim}
\PRset{head=below,logo=below right,smashlogo,ruled=false}
\end{verbatim}
The image is from the \textsf{mwe} package.

Use short paragraphs. Try not to exceed 500 words. Keep to the point
and avoid jargon.

\begin{about}
Some Company was set up in imaginary year $i$.
\end{about}

\end{pressrelease}

\end{document}
